% Options for packages loaded elsewhere
\PassOptionsToPackage{unicode}{hyperref}
\PassOptionsToPackage{hyphens}{url}
\documentclass[
]{article}

% Callouts
\usepackage[most]{tcolorbox}
\newtcbtheorem[auto counter, number within=section]{theorem}{Theorem}{fonttitle=\bfseries, colback=cyan!5!white,colframe=cyan!50!black}{thm}
\newtcbtheorem[auto counter, number within=section]{definition}{Definition}{fonttitle=\bfseries, colback=gray!5!white,colframe=gray!20!black}{def}
\newtcbtheorem[auto counter, number within=section]{lemma}{Lemma}{fonttitle=\bfseries, colback=green!5!white,colframe=green!20!black}{lem}
\newtcbtheorem[auto counter, number within=section]{proposition}{Proposition}{fonttitle=\bfseries, colback=yellow!5!white,colframe=yellow!20!black}{prop}
\newtcbtheorem[auto counter, number within=section]{example}{Example}{fonttitle=\bfseries, colback=purple!7!white,colframe=purple!20!gray}{ex}
\newtcolorbox{info}{title=Info, fonttitle=\bfseries, colback=cyan!5!white,colframe=cyan!30!black}
\newtcolorbox{reference}{title=Reference, fonttitle=\bfseries, colback=gray!5!white,colframe=gray!50!black}
\newtcolorbox{summary}{title=Summary, fonttitle=\bfseries, colback=magenta!5!white,colframe=magenta!30!black}
\newtcolorbox{hint}{title=Hint, fonttitle=\bfseries, colback=green!5!white,colframe=green!30!black}
\newtcolorbox{tip}{title=Tip, fonttitle=\bfseries, colback=green!5!white,colframe=green!30!black}
\newtcolorbox{question}{title=Question, fonttitle=\bfseries, colback=orange!5!white,colframe=green!30!black}
\newtcolorbox{note}{title=Note, fonttitle=\bfseries, colback=cyan!5!white,colframe=cyan!20!black}
\newtcolorbox{attention}{title=Attention, fonttitle=\bfseries, colback=orange!10!white,colframe=orange!75!white}
\newtcolorbox{tldr}{title=TL;DR, fonttitle=\bfseries, colback=cyan!10!white,colframe=cyan!75!white}

\usepackage{xcolor}
\usepackage{amsmath,amssymb}
\setcounter{secnumdepth}{5}
\usepackage{iftex}
\ifPDFTeX
  \usepackage[T1]{fontenc}
  \usepackage[utf8]{inputenc}
  \usepackage{textcomp} % provide euro and other symbols
\else % if luatex or xetex
  \usepackage{unicode-math} % this also loads fontspec
  \defaultfontfeatures{Scale=MatchLowercase}
  \defaultfontfeatures[\rmfamily]{Ligatures=TeX,Scale=1}
\fi
\usepackage{lmodern}
\makeatletter
\@ifundefined{KOMAClassName}{% if non-KOMA class
  \IfFileExists{parskip.sty}{%
    \usepackage{parskip}
  }{% else
    \setlength{\parindent}{0pt}
    \setlength{\parskip}{6pt plus 2pt minus 1pt}}
}{% if KOMA class
  \KOMAoptions{parskip=half}}
\makeatother
\usepackage{color}
\usepackage{fancyvrb}
\newcommand{\VerbBar}{|}
\newcommand{\VERB}{\Verb[commandchars=\\\{\}]}
\DefineVerbatimEnvironment{Highlighting}{Verbatim}{commandchars=\\\{\}}
% Add ',fontsize=\small' for more characters per line
\newenvironment{Shaded}{}{}
\newcommand{\AlertTok}[1]{\textcolor[rgb]{1.00,0.00,0.00}{\textbf{#1}}}
\newcommand{\AnnotationTok}[1]{\textcolor[rgb]{0.38,0.63,0.69}{\textbf{\textit{#1}}}}
\newcommand{\AttributeTok}[1]{\textcolor[rgb]{0.49,0.56,0.16}{#1}}
\newcommand{\BaseNTok}[1]{\textcolor[rgb]{0.25,0.63,0.44}{#1}}
\newcommand{\BuiltInTok}[1]{\textcolor[rgb]{0.00,0.50,0.00}{#1}}
\newcommand{\CharTok}[1]{\textcolor[rgb]{0.25,0.44,0.63}{#1}}
\newcommand{\CommentTok}[1]{\textcolor[rgb]{0.38,0.63,0.69}{\textit{#1}}}
\newcommand{\CommentVarTok}[1]{\textcolor[rgb]{0.38,0.63,0.69}{\textbf{\textit{#1}}}}
\newcommand{\ConstantTok}[1]{\textcolor[rgb]{0.53,0.00,0.00}{#1}}
\newcommand{\ControlFlowTok}[1]{\textcolor[rgb]{0.00,0.44,0.13}{\textbf{#1}}}
\newcommand{\DataTypeTok}[1]{\textcolor[rgb]{0.56,0.13,0.00}{#1}}
\newcommand{\DecValTok}[1]{\textcolor[rgb]{0.25,0.63,0.44}{#1}}
\newcommand{\DocumentationTok}[1]{\textcolor[rgb]{0.73,0.13,0.13}{\textit{#1}}}
\newcommand{\ErrorTok}[1]{\textcolor[rgb]{1.00,0.00,0.00}{\textbf{#1}}}
\newcommand{\ExtensionTok}[1]{#1}
\newcommand{\FloatTok}[1]{\textcolor[rgb]{0.25,0.63,0.44}{#1}}
\newcommand{\FunctionTok}[1]{\textcolor[rgb]{0.02,0.16,0.49}{#1}}
\newcommand{\ImportTok}[1]{\textcolor[rgb]{0.00,0.50,0.00}{\textbf{#1}}}
\newcommand{\InformationTok}[1]{\textcolor[rgb]{0.38,0.63,0.69}{\textbf{\textit{#1}}}}
\newcommand{\KeywordTok}[1]{\textcolor[rgb]{0.00,0.44,0.13}{\textbf{#1}}}
\newcommand{\NormalTok}[1]{#1}
\newcommand{\OperatorTok}[1]{\textcolor[rgb]{0.40,0.40,0.40}{#1}}
\newcommand{\OtherTok}[1]{\textcolor[rgb]{0.00,0.44,0.13}{#1}}
\newcommand{\PreprocessorTok}[1]{\textcolor[rgb]{0.74,0.48,0.00}{#1}}
\newcommand{\RegionMarkerTok}[1]{#1}
\newcommand{\SpecialCharTok}[1]{\textcolor[rgb]{0.25,0.44,0.63}{#1}}
\newcommand{\SpecialStringTok}[1]{\textcolor[rgb]{0.73,0.40,0.53}{#1}}
\newcommand{\StringTok}[1]{\textcolor[rgb]{0.25,0.44,0.63}{#1}}
\newcommand{\VariableTok}[1]{\textcolor[rgb]{0.10,0.09,0.49}{#1}}
\newcommand{\VerbatimStringTok}[1]{\textcolor[rgb]{0.25,0.44,0.63}{#1}}
\newcommand{\WarningTok}[1]{\textcolor[rgb]{0.38,0.63,0.69}{\textbf{\textit{#1}}}}
\usepackage{longtable,booktabs,array}
\newcounter{none} % for unnumbered tables
\usepackage{calc} % for calculating minipage widths
% Correct order of tables after \paragraph or \subparagraph
\usepackage{etoolbox}
\makeatletter
\patchcmd\longtable{\par}{\if@noskipsec\mbox{}\fi\par}{}{}
\makeatother
% Allow footnotes in longtable head/foot
\IfFileExists{footnotehyper.sty}{\usepackage{footnotehyper}}{\usepackage{footnote}}
\makesavenoteenv{longtable}
\ifLuaTeX
  \usepackage{luacolor}
  \usepackage[soul]{lua-ul}
\else
  \usepackage{soul}
\fi
\setlength{\emergencystretch}{3em} % prevent overfull lines
\providecommand{\tightlist}{%
  \setlength{\itemsep}{0pt}\setlength{\parskip}{0pt}}
\usepackage{fontsetup}
\usepackage[a4paper,left=20mm,top=20mm]{geometry}
\usepackage{bookmark}
\IfFileExists{xurl.sty}{\usepackage{xurl}}{} % add URL line breaks if available
\urlstyle{same}
\hypersetup{
  hidelinks,
  pdfcreator={LaTeX via pandoc}}

\author{}
\date{}

\begin{document}

\section{Plan: Clean Rapier-based sim-dev rewrite
(v2-only)}\label{plan-clean-rapier-based-sim-dev-rewrite-v2-only}

\subsection{Status}\label{status}

{\def\LTcaptype{none} % do not increment counter
\begin{longtable}[]{@{}
  >{\raggedright\arraybackslash}p{(\linewidth - 4\tabcolsep) * \real{0.3333}}
  >{\raggedright\arraybackslash}p{(\linewidth - 4\tabcolsep) * \real{0.3333}}
  >{\raggedright\arraybackslash}p{(\linewidth - 4\tabcolsep) * \real{0.3333}}@{}}
\toprule\noalign{}
\begin{minipage}[b]{\linewidth}\raggedright
Phase
\end{minipage} & \begin{minipage}[b]{\linewidth}\raggedright
State
\end{minipage} & \begin{minipage}[b]{\linewidth}\raggedright
Notes
\end{minipage} \\
\midrule\noalign{}
\endhead
\bottomrule\noalign{}
\endlastfoot
1 --- Project skeleton + scene & ✅ done & \\
2 --- Robot loader (no physics) & ✅ done & \\
3 --- Static physics + drop test & ✅ done & \\
4 --- Map loading (JSON parity) & ✅ done & model loading async; trimesh
colliders for \texttt{immovable} \\
5 --- Custom raycast wheel force model & ✅ done & replaced
\texttt{DynamicRayCastVehicleController}; drift / turn fixes landed
(\texttt{e52aca9}) \\
5b --- UI scaffolding (camera modes, splash, presets, text labels,
remember-stage, runtime controls) & ✅ done & \\
6 --- Stage object collider visualization & ✅ done & wireframes for
floor / cube / cylinder / model \\
7 --- Debug menu (lil-gui) & ✅ done & all folders wired; most params
mutate live (mass/radius require body recreation) \\
8 --- Ramps and behavior validation & ✅ done & holds on all ramps
(5°--40°); no slip; descent controlled; edge drop clean \\
9 --- Replace \texttt{sim-dev/} with \texttt{sim-v2/} & ⬜ not started &
needs \texttt{React.StrictMode} re-enable first \\
\end{longtable}
}

\subsection{Context}\label{context}

The current \texttt{sim-dev/} drives the robot with a kinematic-ish
approach: 2 sphere wheel colliders attached to the chassis body,
velocities set directly via \texttt{setLinvel}/\texttt{setAngvel}, plus
an uprightness gate and wheel raycast probes. This works on flat ground
but \textbf{produces} wrong wheel positions, no real suspension, no
slip, no holding-on-ramps behavior. The plan: replace it with a small
\textbf{custom raycast wheel force model} --- per-wheel downward raycast
→ suspension spring force + per-wheel longitudinal traction
(engine/brake) + side-friction force, all applied through
\texttt{RigidBody.applyImpulseAtPoint} on the chassis. (Earlier draft
used Rapier's \texttt{DynamicRayCastVehicleController}; that was dropped
--- see Phase 5 below.)

The current \texttt{sim-dev/} also depends on
\texttt{front-end/src/components/js-simulator/} (the V1 OBJ loader,
scene/camera/lights, stage JSON loader). The new tree will not. Once the
new tree drives v2 cleanly on flat ground + ramps, it replaces
\texttt{sim-dev/}.

The intended outcome is a small, owned, easy-to-tune simulator that gets
wheel mechanics \emph{right} --- correct positions, real suspension,
real friction --- so the rest of the GSoC project (sensors, programs,
multi-robot) sits on a trustworthy foundation.

\subsection{Scope (v2 only, no V1, no js-simulator
import)}\label{scope-v2-only-no-v1-no-js-simulator-import}

\begin{itemize}
\tightlist
\item
  One robot: v2, loaded from
  \texttt{sim-v2/public/js-simulator/models/robots/v2/*.stl} so this
  tree owns the assets it serves.
\item
  Stages: a flat plane stub (phases 1--3) replaced by JSON stages with
  parity to the V1 js-simulator schema (phase 4); ramp set added in
  phase 8.
\item
  Two driving wheels (raycast) + one caster (passive free-roll sphere
  collider on chassis --- see ``Caster decision'' below).
\item
  Debug GUI via \texttt{lil-gui} (small, well-known, dat.gui-style
  folders/sliders/buttons).
\end{itemize}

\subsection{New directory}\label{new-directory}

\texttt{sim-v2/} at repo root, sibling to \texttt{sim-dev/}. Same Vite +
React + TS stack so the eventual swap is mechanical. No imports from
\texttt{front-end/src/components/js-simulator/}.

\textbf{Asset strategy:} STL parts for v2 are loaded from the existing
public URL path (\texttt{/js-simulator/models/robots/v2/*.stl}) --- they
stay where they are, the new tree just fetches them. The wheel mesh
(\texttt{wheel.obj} + \texttt{wheel.mtl}) is \textbf{copied} into
\texttt{sim-v2/public/wheel/} so the new tree owns its own wheel asset
and doesn't import from the V1 robot folder.

\subsection{Caster decision (load-bearing, picked here so Phase 3
doesn't
regress)}\label{caster-decision-load-bearing-picked-here-so-phase-3-doesnt-regress}

\texttt{DynamicRayCastVehicleController} is car-shaped: N steered/driven
raycast wheels, no concept of a passive caster. Three options were
considered:

\begin{itemize}
\tightlist
\item
  \textbf{(A)} 2 raycast wheels for drive + caster as a free
  \texttt{Ball} collider on the chassis. Caster contributes weight,
  contact, and stability the way the real robot's caster does. Selected.
\item
  \begin{enumerate}
  \def\labelenumi{(\Alph{enumi})}
  \setcounter{enumi}{1}
  \tightlist
  \item
    3 raycast wheels (caster as a third raycast wheel,
    non-driven/non-steered). Cheaper but loses the rolling-ball physics
    of a real caster and is a hack against the controller's intent.
  \end{enumerate}
\item
  \begin{enumerate}
  \def\labelenumi{(\Alph{enumi})}
  \setcounter{enumi}{2}
  \tightlist
  \item
    Lift caster out of physics entirely (visual only). Wrong --- the
    caster carries a third of the weight; ignoring it changes balance.
  \end{enumerate}
\end{itemize}

Going with (A). If (A) misbehaves on ramps (caster ball sticking,
popping), fall back to (B).

\subsection{Phase ordering (each phase ends with a verifiable
checkpoint)}\label{phase-ordering-each-phase-ends-with-a-verifiable-checkpoint}

\subsubsection{Phase 1 --- Project skeleton + scene (no physics, no
robot)}\label{phase-1-project-skeleton-scene-no-physics-no-robot}

\begin{enumerate}
\def\labelenumi{\arabic{enumi}.}
\tightlist
\item
  Scaffold \texttt{sim-v2/} (package.json, tsconfig, vite.config,
  index.html, \texttt{src/main.tsx}).
\item
  Bootstrap Three.js: renderer, scene, perspective camera,
  OrbitControls, ambient + directional light, ground plane
  (PlaneGeometry, 20×20m, gridhelper on top), RAF render loop.
\item
  Top-right \textbf{orientation gizmo}: a second small
  \texttt{WebGLRenderer} viewport (or a small scene rendered into a
  corner via \texttt{setViewport}) showing an \texttt{AxesHelper} whose
  rotation tracks the robot's chassis rotation. World axes can be a
  separate static helper for reference.
\item
  Window resize + DPR handling.
\end{enumerate}

\textbf{Verify:} open the page, see grid + lighting + orbit camera, see
two axis triads (world reference + corner gizmo) responding to the
camera.

\subsubsection{Phase 2 --- Load v2 robot, correct orientation, no
physics}\label{phase-2-load-v2-robot-correct-orientation-no-physics}

\begin{enumerate}
\def\labelenumi{\arabic{enumi}.}
\tightlist
\item
  STL loader (\texttt{STLLoader} from
  \texttt{three/examples/jsm/loaders/STLLoader}).
\item
  Reimplement the v2 part assembler. Mirror the structure of
  \texttt{sim-dev/src/robot/robotV2.ts} for known-good values:

  \begin{itemize}
  \tightlist
  \item
    Outer pivot \texttt{Group} with \texttt{rotation.x\ =\ -π/2} (Fusion
    Z-up → world Y-up).
  \item
    Inner parts at \texttt{PART\_DEFAULTS} mm offsets, scaled by 0.001.
  \item
    Right-fender mirror trick (right\_fender.stl is broken --- mirror
    left).
  \item
    \texttt{front\_base.rotation.z\ =\ π} (face forward).
  \item
    Color map per part.
  \end{itemize}
\item
  Wheels: load \texttt{caster.obj} once (caster) + a wheel STL/OBJ
  (TODO: confirm asset --- see Q2 below). Place left/right wheels at
  fender-center mount points using the existing constants.
\item
  Caster mesh placed at its mount point.
\end{enumerate}

\textbf{Verify:} robot stands on grid plane, body upright, wheels
visually at the fender mounts, caster underneath at the right spot,
orientation gizmo's axes match what the robot looks like (robot's
``forward'' matches gizmo's forward).

\subsubsection{Phase 3 --- Static physics + chassis drop test (no
vehicle controller
yet)}\label{phase-3-static-physics-chassis-drop-test-no-vehicle-controller-yet}

\begin{enumerate}
\def\labelenumi{\arabic{enumi}.}
\tightlist
\item
  Init Rapier (\texttt{@dimforge/rapier3d-compat}):
  \texttt{await\ RAPIER.init()}, world with gravity
  \texttt{(0,-9.81,0)}.
\item
  Static ground: cuboid collider at y=0.
\item
  Chassis rigid body (dynamic). For collision shape use the v2
  \texttt{\_coacd} convex hulls (compound \texttt{ConvexHull} colliders)
  --- this asset set is already in the repo. Mass set explicitly (target
  \textasciitilde1.5--2.5 kg total --- confirm in Q3).
\item
  Caster: \texttt{Ball} collider on the same body, radius from caster
  mesh bbox.
\item
  Mesh sync: each frame, copy \texttt{body.translation()} and
  \texttt{body.rotation()} to the visual pivot, applying the body-local
  visual offset rotated by the body quaternion (the fix from
  \texttt{FINDINGS.md}).
\item
  Stepping: \texttt{world.step()} at fixed dt (1/60), accumulator
  pattern in the RAF loop.
\end{enumerate}

\textbf{Verify:} spawn robot 0.5m above ground, watch it fall and rest
stably. Body upright, caster contacting, no jitter, mesh stays glued to
body.

\subsubsection{Phase 4 --- Map loading (JSON stage
parity)}\label{phase-4-map-loading-json-stage-parity}

Goal: load existing js-simulator stages without modification. Stage
files are bundled from \texttt{sim-v2/src/stages/data/*.json} and
referenced assets live under \texttt{sim-v2/public/js-simulator/}. The
schema and asset URLs match V1, so the same JSON shape drives both
simulators.

\textbf{Schema observed in the V1 stages} (no formal spec, derived from
reading \texttt{stage\_white\_rect.json}, \texttt{stage\_ramp.json},
\texttt{stage\_object.json}):

\begin{itemize}
\tightlist
\item
  \texttt{floor} --- \texttt{dimensions:\ {[}w,\ h{]}} (per-tile;
  texture tiles via \texttt{repeat:\ {[}u,\ v{]}}),
  \texttt{material.color} (string or decimal RGB), optional
  \texttt{texture} URL, optional \texttt{offset}.
\item
  \texttt{cube} --- \texttt{dimensions:\ {[}w,\ h,\ d{]}},
  \texttt{material.color}, \texttt{position}, optional \texttt{mass}
  (kg; absent / 0 / \texttt{immovable:\ true} → static), optional
  \texttt{castShadow}.
\item
  \texttt{cylinder} ---
  \texttt{dimensions:\ {[}radiusTop,\ radiusBottom,\ height,\ radialSegments{]}},
  \texttt{material.color}, \texttt{position}, optional \texttt{mass},
  optional \texttt{immovable}.
\item
  \texttt{model} --- \texttt{filename} (absolute URL to an OBJ),
  \texttt{position}, optional \texttt{scale} (uniform), optional
  \texttt{orientation:\ {[}rx,\ ry,\ rz{]}} (radians, XYZ Euler),
  optional \texttt{color}, \texttt{castShadow}, \texttt{immovable}.
\item
  \texttt{fossbot} --- robot spawn marker. \texttt{position} +
  \texttt{orientation}. Exactly one per stage.
\end{itemize}

Common fields: \texttt{name} (debug string),
\texttt{position:\ {[}x,\ y,\ z{]}} in meters Y-up, \texttt{orientation}
in radians.

\begin{enumerate}
\def\labelenumi{\arabic{enumi}.}
\tightlist
\item
  \textbf{Loader} --- \texttt{sim-v2/src/stages/loader.ts} exports
  \texttt{loadStage(name:\ string,\ scene,\ world):\ Promise\textless{}StageHandle\textgreater{}}.
  Fetches \texttt{/stages/\$\{name\}.json}, parses, dispatches each
  entry to a builder.
\item
  \textbf{Builders} --- one per object type, each returns
  \texttt{\{\ visuals:\ Object3D{[}{]},\ colliders:\ Collider{[}{]}\ \}}:

  \begin{itemize}
  \tightlist
  \item
    \texttt{buildFloor} --- \texttt{THREE.PlaneGeometry} rotated to lie
    on Y=0; large static \texttt{cuboid(w/2,\ 0.005,\ h/2)} collider
    replacing the Phase 3 stub ground; loads texture via
    \texttt{THREE.TextureLoader} with \texttt{repeat}/\texttt{offset} if
    present.
  \item
    \texttt{buildCube} / \texttt{buildCylinder} --- Three primitive +
    Rapier \texttt{cuboid} / \texttt{cylinder} collider; dynamic body if
    \texttt{mass\ \textgreater{}\ 0} and not \texttt{immovable}, else
    static (\texttt{fixed}).
  \item
    \texttt{buildModel} --- \texttt{OBJLoader} (already in
    \texttt{three/examples/jsm}); for \texttt{immovable} items use
    \texttt{RAPIER.ColliderDesc.trimesh(vertices,\ indices)} from the
    loaded geometry; for dynamic ones fall back to a single oriented
    bounding box (good enough for v2's needs --- full convex decomp is
    the kind of thing CoACD would give us, deferred).
  \item
    \texttt{buildFossbot} --- sets the chassis body's
    translation/rotation to the spawn pose; does NOT recreate the body.
  \end{itemize}
\item
  \textbf{Color parsing} --- accept both
  \texttt{\textquotesingle{}orange\textquotesingle{}} and decimal
  numbers (\texttt{65280}). One helper
  \texttt{parseColor(value):\ THREE.ColorRepresentation}.
\item
  \textbf{Asset paths} --- JSON references absolute URLs like
  \texttt{/js-simulator/models/static/Ramp.obj} and
  \texttt{/js-simulator/textures/white\_rect.png};
  \texttt{sim-v2/public/js-simulator/} owns those files locally so the
  simulator no longer depends on \texttt{front-end/public}.

  \begin{itemize}
  \tightlist
  \item
    Start with (a). Fall back to (b) only if Vite proxying gets in the
    way.
  \end{itemize}
\item
  \textbf{Stage handle} --- \texttt{StageHandle} exposes
  \texttt{dispose()} that removes meshes from the scene and removes
  static colliders from the world (track them for clean teardown).
  Required for the lil-gui stage switcher in the next phase.
\item
  \textbf{Stage registry} --- \texttt{sim-v2/src/stages/index.ts}
  enumerates known stages by listing the names already shipped under
  \texttt{public/stages/}. Used by the GUI dropdown in Phase 6.
\end{enumerate}

\textbf{Verify:} - \texttt{stage\_white\_rect}, \texttt{stage\_ramp},
\texttt{stage\_object} all load without error and look visually
equivalent to V1. - Robot spawns at the position the JSON says, not the
hardcoded Phase 3 spawn. - Static cubes/cylinders/ramps are immovable;
dynamic ones (mass \textgreater{} 0) get pushed when the robot bumps
them. - Switching from one stage to another via a temporary keyboard
shortcut leaves no leaked colliders (check with
\texttt{world.bodies.len()} / collider count).

\subsubsection{Phase 5 --- Custom raycast wheel force model (replaces
vehicle
controller)}\label{phase-5-custom-raycast-wheel-force-model-replaces-vehicle-controller}

Goal: replace the static cylinder wheel colliders with a per-wheel
raycast force model that gives real suspension, real traction, and clean
ramp behavior on a 2-wheel differential-drive chassis. Caster stays a
passive \texttt{Ball} collider; only the two drive wheels go through the
model.

\textbf{Why not \texttt{DynamicRayCastVehicleController}?} It was the
initial design (and the entire previous draft of this phase). Dropped
during ramp testing because:

\begin{itemize}
\tightlist
\item
  The controller is car-shaped: it assumes N steered/driven raycast
  wheels. v2 is \textbf{two-wheeled with fixed (non-steering,
  non-rotating) wheels} --- most of the controller's machinery (steering
  angle, axle Cs, wheel rotation tracking) is dead weight or actively
  fights the differential-drive model.
\item
  The controller writes chassis velocities directly each
  \texttt{updateVehicle()}, which fights any external torque or impulse
  we apply (tried this for in-place rotation in an earlier session ---
  see ``Option 1 torque'' history; the controller's lateral-friction
  term anchored rotation to the wheel-axle line, not the chassis CoM).
\item
  On inclines the velocity-write made it hard to add a slope-hold term
  without arguing with the controller.
\end{itemize}

The custom model applies forces (not velocities) at each wheel contact
point, leaving the chassis as a normal Rapier dynamic body integrated by
\texttt{world.step()}.

\textbf{Pre-step --- drop the wheel cylinder colliders.} The original
\texttt{ROBOT\_COLLIDERS} left/right cylinders are still filtered out of
the chassis build (via
\texttt{createRobotBody(\{\ skipDriveWheels:\ true\ \})}). Two solid
colliders riding the ground while we also raycast = contact thrash.

\textbf{Per-frame model (in \texttt{physics/vehicle.ts}).} For each
wheel \texttt{i\ ∈\ \{LEFT,\ RIGHT\}}:

\begin{enumerate}
\def\labelenumi{\arabic{enumi}.}
\tightlist
\item
  Compute the wheel attachment point in world space (chassis-local
  connection point transformed by chassis pose).
\item
  Cast a ray downward (chassis-local \texttt{(0,-1,0)} rotated to world)
  from the attachment, length
  \texttt{suspensionRestLength\ +\ maxSuspensionTravel\ +\ wheelRadius}.
\item
  If the ray hits the floor:

  \begin{itemize}
  \tightlist
  \item
    \textbf{Suspension force}: spring along the suspension axis,
    magnitude \texttt{k\ *\ compression\ -\ c\ *\ compressionVelocity}.
    Applied at the wheel contact point on the chassis body.
  \item
    \textbf{Longitudinal force} (engine / brake):
    \texttt{motorForce\ *\ input{[}i{]}} along chassis-forward at the
    contact point. \texttt{input{[}i{]}\ ∈\ {[}-1,\ 1{]}} from
    \texttt{setDrive(left,\ right)}. Brake is applied as opposing force
    when input is zero or sign-flipped.
  \item
    \textbf{Side-friction force}: opposite to the lateral component of
    the contact-point velocity, scaled by \texttt{sideFriction}. Keeps
    the chassis from sliding sideways while turning.
  \item
    \textbf{Slope-hold}: an extra anti-rollback force
    \texttt{chassis.mass()\ *\ g\ *\ sin(slope)\ *\ slopeFactor} along
    chassis-forward, scaled to the wheel's contact normal. Lets the
    robot brake-hold on a steep ramp without engine input.
  \end{itemize}
\item
  If the ray misses, the wheel is airborne --- no force applied this
  frame.
\end{enumerate}

\textbf{Tunable settings (live in the GUI; defaults in
\texttt{createDefaultVehicleSettings(wheelRadius)}).}

\begin{itemize}
\tightlist
\item
  \texttt{motorForce} = 12 (Newtons at full throttle, per wheel).
\item
  \texttt{brakeStrength} --- opposing force when input is zero / brake
  is held.
\item
  \texttt{suspensionStiffness}, \texttt{suspensionDamping},
  \texttt{suspensionRestLength}, \texttt{maxSuspensionTravel}.
\item
  \texttt{sideFriction} --- side-friction stiffness for in-place
  stability.
\item
  \texttt{slopeFactor} --- slope-hold gain (0 disables).
\item
  \texttt{wheelRadius} is sourced from the loaded wheel mesh, not
  hardcoded.
\end{itemize}

\textbf{API.}

\begin{Shaded}
\begin{Highlighting}[]
\FunctionTok{createVehicle}\NormalTok{(world}\OperatorTok{,}\NormalTok{ chassis)}\OperatorTok{:}\NormalTok{ VehicleHandle}
\NormalTok{VehicleHandle}\OperatorTok{.}\FunctionTok{setDrive}\NormalTok{(left}\OperatorTok{:} \DataTypeTok{number}\OperatorTok{,}\NormalTok{ right}\OperatorTok{:} \DataTypeTok{number}\NormalTok{)  }\CommentTok{// each in {-}1..1}
\NormalTok{VehicleHandle}\OperatorTok{.}\FunctionTok{update}\NormalTok{(dt)                             }\CommentTok{// call before world.step()}
\NormalTok{VehicleHandle}\OperatorTok{.}\AttributeTok{telemetry}\OperatorTok{:}\NormalTok{ VehicleTelemetry            }\CommentTok{// per{-}wheel contact, susp, slip, force}
\end{Highlighting}
\end{Shaded}

\texttt{setDrive} only wakes the chassis when input is non-zero (lets
the body sleep when idle, fixing the earlier idle-drift).

\textbf{Differential drive mapping (in \texttt{App.tsx} keyboard
handler).}

\begin{itemize}
\tightlist
\item
  W/S → both wheels ±1.
\item
  A/D → opposite signs for in-place rotation, or asymmetric magnitudes
  when combined with W/S for skid-steer.
\item
  Space → idle/brake (\texttt{setDrive(0,\ 0)} and chassis sleeps when
  velocity decays).
\end{itemize}

\textbf{Visual wheel sync.} Wheels are non-steering (the chassis is
two-wheeled differential-drive --- no Ackermann, no yaw-axis rotation on
the wheels themselves), but they do \textbf{spin visually} around their
body-local X axle and bob with the suspension:

\begin{itemize}
\tightlist
\item
  Roll:
  \texttt{wheelSpin{[}i{]}\ +=\ (longitudinalContactVelocity{[}i{]}\ /\ wheelRadius)\ *\ dt}
  while the wheel is touching ground; falls back to
  \texttt{input\ *\ freeSpinSpeed\ *\ dt} when airborne so the wheel
  still spins under throttle.
  \texttt{visualWheels{[}i{]}.rotation.x\ =\ wheelSpin{[}i{]}}.
\item
  Suspension bob:
  \texttt{visualWheels{[}i{]}.position.y\ =\ baseY\ +\ (suspensionRestLength\ −\ currentSuspensionLength)}
  so the mesh dips into the fender as the spring compresses.
\end{itemize}

The wheel meshes stay parented under \texttt{robot.root} (the chassis
frame) --- only the local position.y and rotation.x are touched per
frame.

\textbf{Verify:} - WASD drives the robot. - Forward/reverse symmetric. -
A/D rotates roughly in place around the chassis center; D = clockwise
from above. - No drift on flat ground when no input. - Robot brake-holds
on a 30° ramp (slope-hold factor active). - Suspension visibly works
(telemetry reports compression on bumps).

\subsubsection{Phase 5b --- UI scaffolding (parallel to Phase
5)}\label{phase-5b-ui-scaffolding-parallel-to-phase-5}

Goal: make the simulator usable for hand-tuning during ramp work. These
pieces all landed in the May 5--6 ramp-testing burst and support Phase
5/Phase 8 iteration. They are visual / UX scaffolding, not physics ---
kept in their own phase so Phase 7 (Debug menu) stays focused on lil-gui
parameter binding.

\begin{enumerate}
\def\labelenumi{\arabic{enumi}.}
\tightlist
\item
  \textbf{Camera modes} (\texttt{src/ui/cameraControls.ts},
  \texttt{src/ui/cameraTypes.ts}). A view-mode button next to the Debug
  GUI cycles \textbf{Orbit → Follow → Top}:

  \begin{itemize}
  \tightlist
  \item
    Orbit: the existing OrbitControls camera; pose is preserved when
    leaving and restored when returning.
  \item
    Follow: third-person chase camera tracking behind the robot.
  \item
    Top: fixed eagle-eye for surveying the whole stage.
  \end{itemize}
\item
  \textbf{Splash screen} (\texttt{src/ui/splashScreen.ts}). Startup
  overlay reporting load progress (world init, stage load, robot load,
  vehicle ready). Visibility persisted in the World folder of the Debug
  GUI.
\item
  \textbf{Movement presets} (\texttt{src/ui/movementPresets.ts}).
  Physics-driven keyboard shortcuts for canned moves (e.g.~spawn, drive
  forward N seconds, pivot 90°). Used to repro tuning scenarios
  reliably.
\item
  \textbf{\texttt{text} stage entry type} (\texttt{buildText} in
  \texttt{stages/builders.ts}). Canvas-backed sprite label, optional
  \texttt{onFloor} flag for ground-aligned labels. Used in
  \texttt{stage\_ramp.json} to call out which lane is which fixture.
\item
  \textbf{``Remember last stage'' toggle}
  (\texttt{src/debug/...stageFolder\ +\ localStorage}). When on, the
  last selected stage is restored on reload.
\item
  \textbf{Reset-to-spawn} action exposed in the Stage folder of the
  Debug GUI; same path used after a stage swap.
\item
  \textbf{Runtime controls} wired in \texttt{App.tsx} and the World
  folder: pause/step/time-scale, telemetry overlay visibility, stage
  collider wireframe visibility (preserved across stage swaps).
\end{enumerate}

\textbf{Verify:} - Switching camera modes is instant; orbit pose
round-trips when toggling away and back. - Splash screen disappears once
everything is ready; reappears on hard reload. - Stage swap →
reset-to-spawn → drive: clean, no leaked state. - ``Remember last
stage'' survives a page reload. - Pause + step-once advances physics by
exactly one fixed step.

\subsubsection{Phase 6 --- Stage object collider
visualization}\label{phase-6-stage-object-collider-visualization}

Goal: extend the existing ``Show Colliders'' toggle (which already
wireframes the robot's primitive colliders) to also wireframe stage
object colliders, so when something looks off in physics it's obvious
whether the visual or the collider is wrong.

\begin{enumerate}
\def\labelenumi{\arabic{enumi}.}
\tightlist
\item
  \textbf{Builders return collider metadata.} Each stage builder already
  returns \texttt{\{\ object,\ collider\ \}}. Add a sibling debug
  wireframe (matching the collider's shape and transform) and a small
  marker so the loader can group them.
\item
  \textbf{Stage builds a hidden colliders group.} \texttt{loader.ts}
  collects every debug wireframe into a single \texttt{THREE.Group}
  named \texttt{stage\_colliders}, parented to the scene root,
  \texttt{visible\ =\ false} by default.
\item
  \textbf{Toggle wires through the existing GUI.} Reuse the Debug GUI's
  ``Show Colliders'' boolean --- flipping it toggles both the robot
  colliders group AND the active stage's colliders group.
\item
  \textbf{Trimesh models}: bake the same vertices/indices used for the
  Rapier \texttt{trimesh} collider into a \texttt{THREE.LineSegments}
  (\texttt{EdgesGeometry} over a \texttt{BufferGeometry}) so wireframes
  match the physics exactly, including baked transforms.
\item
  \textbf{Stage swap cleanup.} \texttt{StageHandle.dispose()} removes
  the colliders group along with the visual objects.
\end{enumerate}

\textbf{Verify:} - Toggling ``Show Colliders'' with
\texttt{stage\_object} shows magenta cube/cylinder wireframes co-located
with the visual cubes/cylinders. - Toggling on \texttt{stage\_ramp}
shows the trimesh wireframe of \texttt{Ramp.obj} aligning with the
visible mesh. - Toggling off restores the clean view. - Stage swap
leaves no orphan wireframes in the scene.

\subsubsection{Phase 7 --- Debug menu (lil-gui) and reload
affordances}\label{phase-7-debug-menu-lil-gui-and-reload-affordances}

\begin{enumerate}
\def\labelenumi{\arabic{enumi}.}
\tightlist
\item
  Install \texttt{lil-gui}. Mount one panel, top-left, collapsible.
\item
  Folders + bindings:

  \begin{itemize}
  \tightlist
  \item
    \textbf{World}: gravity, time scale, paused, physics wireframes
    (\texttt{world.debugRender}), step-once button.
  \item
    \textbf{Robot}: mass, body linear/angular damping, reset-to-spawn
    button, log-state button. Manual movement along axes.
  \item
    \textbf{Wheels} (per side or shared): radius, suspension stiffness,
    max travel, friction slip, side friction, compression, relaxation,
    engine force max, brake force max.
  \item
    \textbf{Drive}: forward speed cap, turn force scale, deadzone.
  \item
    \textbf{Stage}: dropdown populated from the Phase 4 stage registry
    (e.g.~\texttt{stage\_white\_rect}, \texttt{stage\_ramp},
    \texttt{stage\_object}, plus the new ramp set added in Phase 7).
    Selecting calls the loader's \texttt{dispose()} on the current stage
    and \texttt{loadStage()} on the new one.
  \end{itemize}
\item
  \textbf{Console output panel}: small DOM overlay (not lil-gui)
  printing every N frames: position, velocity, yaw, per-wheel suspension
  length, per-wheel ground contact, per-wheel slip. A ``log now'' button
  dumps the full state to \texttt{console.log} for copy/paste.
\end{enumerate}

\textbf{Verify:} every parameter changes behavior live (no rebuild
needed except mass/radius which need rigid-body recreation). Stage
reload swaps geometry without leaking colliders. Logs are useful, not
spammy.

\subsubsection{Phase 8 --- Ramps and behavior
validation}\label{phase-8-ramps-and-behavior-validation}

\begin{enumerate}
\def\labelenumi{\arabic{enumi}.}
\tightlist
\item
  \textbf{One combined stage},
  \texttt{sim-v2/public/stages/stage\_ramp.json}, with side-by-side
  fixtures so we can tune by driving between lanes without reloading.
  Each fixture is either a tilted \texttt{cube} collider (oriented cube
  support added in Phase 4) or a \texttt{model} referencing a ramp OBJ.
  Floor labels via the new \texttt{text} entry type call out which lane
  is which.
\item
  Fixtures in the combined stage:

  \begin{itemize}
  \tightlist
  \item
    15° short ramp (1 m).
  \item
    15° long ramp (4 m) --- used as an up-flat-down bridge lane.
  \item
    30° steep ramp (2 m).
  \item
    Mixed bump lane: short + long + a small bump for suspension travel.
  \end{itemize}
\item
  Tune the custom raycast model (Phase 5) until:

  \begin{itemize}
  \tightlist
  \item
    \textbf{Climb}: robot climbs the 15° short and long ramps at full
    throttle without stalling.
  \item
    \textbf{Hold}: with throttle released, the slope-hold term keeps the
    robot stationary on the 30° ramp.
  \item
    \textbf{Slip}: with very low \texttt{sideFriction}, robot visibly
    slides sideways on a side-tilted incline --- sanity check.
  \item
    \textbf{Descent}: robot can drive down the 30° ramp controllably
    with brake input; no flipping.
  \item
    \textbf{Edge drop}: robot driving off the side of a ramp drops
    cleanly without sticking.
  \end{itemize}
\end{enumerate}

\textbf{Verify:} record success of each behavior in \texttt{PROGRESS.md}
(one entry per tuning session is fine). Floor labels mean we can
screenshot a single stage to document everything.

\subsubsection{\texorpdfstring{Phase 9 --- Replace
\texttt{sim-dev/}}{Phase 9 --- Replace sim-dev/}}\label{phase-9-replace-sim-dev}

Only after Phase 8 lands cleanly:

\begin{enumerate}
\def\labelenumi{\arabic{enumi}.}
\tightlist
\item
  \texttt{git\ mv\ sim-dev\ sim-dev-old} (or delete after confirming).
\item
  \texttt{git\ mv\ sim-v2\ sim-dev}.
\item
  Update any root-level scripts/links that pointed at \texttt{sim-dev}.
\item
  PROGRESS.md entry.
\end{enumerate}

\subsection{Critical files to mirror values from (read-only references,
not
imports)}\label{critical-files-to-mirror-values-from-read-only-references-not-imports}

\begin{itemize}
\tightlist
\item
  \texttt{sim-dev/src/robot/robotV2.ts} --- \texttt{PART\_DEFAULTS},
  \texttt{PART\_ADJUST}, \texttt{PART\_COLOR}, \texttt{WHEEL\_DEFAULTS},
  fender-center math, axis-correction pivot. Phase 2 reimplements this
  from scratch but copies the \emph{numbers}.
\item
  \texttt{sim-dev/src/physics/world.ts} --- Rapier init pattern
  (\texttt{RAPIER.init()} then \texttt{new\ RAPIER.World(gravity)}).
\item
  \texttt{sim-dev/src/physics/robotBody.ts} --- body-local visual offset
  rotated by body quaternion on sync. Phase 3 reuses this idea.
\item
  \texttt{sim-v2/public/js-simulator/models/robots/v2/*.stl} and
  \texttt{*\_coacd.stl} --- local assets served by the v2 app.
\end{itemize}

\subsection{Verification summary
(end-to-end)}\label{verification-summary-end-to-end}

After all phases, the new tree should pass these by hand:

\begin{enumerate}
\def\labelenumi{\arabic{enumi}.}
\tightlist
\item
  Page loads, robot is upright and centered on a flat plane, gizmo shows
  correct orientation.
\item
  Drop test (toggle physics on): robot settles cleanly on its two wheels
  + caster, no penetration, no jitter.
\item
  WASD drives: forward, reverse, skid-steer turn --- all symmetric and
  predictable.
\item
  Suspension visibly works (wheel meshes compress over bumps).
\item
  Climbs short and long 15° ramps under full throttle.
\item
  Holds on 30° ramp with brake.
\item
  Slip parameter visibly affects traction.
\item
  Edge drop off ramp doesn't get stuck.
\item
  Stage reload from GUI works without leaking colliders.
\item
  Restart-the-robot button respawns at origin without leftover state.
\end{enumerate}

\subsection{Tooling / workflow notes (not phase
work)}\label{tooling-workflow-notes-not-phase-work}

\begin{itemize}
\tightlist
\item
  \textbf{Post-commit hash hook} (\texttt{.githooks/post-commit}, May
  6). After every commit, the hook replaces the active \texttt{PENDING}
  placeholder in \texttt{PROGRESS.md} with the real commit short hash,
  stages the update, and creates a follow-up
  \texttt{chore:\ progress\ commit\ hash} commit. Per-worktree opt-in;
  documented in \texttt{AGENTS.md}.
\item
  \textbf{\texttt{AGENTS.md} replaces \texttt{CLAUDE.md}} (May 6).
  Repository operating rules now live in \texttt{AGENTS.md} (commit
  behavior, progress-log handling, hook automation, simulator-source
  boundaries). The old \texttt{CLAUDE.md} was deleted as superseded.
\end{itemize}

\subsection{Decisions (confirmed)}\label{decisions-confirmed}

\begin{itemize}
\tightlist
\item
  Directory: \texttt{sim-v2/}.
\item
  Wheel mesh: copy \texttt{wheel.obj} + \texttt{wheel.mtl} into
  \texttt{sim-v2/public/wheel/}.
\item
  Caster: free \texttt{Ball} collider on chassis (option A).
\item
  Debug GUI: \texttt{lil-gui}.
\end{itemize}

\subsection{Beyond the plan}\label{beyond-the-plan}

Features and utilities built that were not listed in the original phase
scope:

\begin{itemize}
\tightlist
\item
  \textbf{\texttt{positionPresets} / \texttt{positionStore}} ---
  save/restore robot poses per stage; persisted to
  \texttt{localStorage}.
\item
  \textbf{\texttt{collidersTunerFolder}} --- per-collider live tuning
  (size, position, rotation, friction, density, restitution) with
  dump-to-clipboard and reset-to-defaults.
\item
  \textbf{\texttt{visualTunerFolder}} --- fender and wheel position
  tuning with CAD-mm units; mirror toggle; dump emits paste-ready
  constants for \texttt{v2.ts}.
\item
  \textbf{\texttt{checkSuspensionHealth}} (\texttt{util/suspension.ts})
  --- static sanity check for suspension stiffness vs robot weight; logs
  critical/soft/ok verdict.
\item
  \textbf{Telemetry overlay copy} --- \texttt{copyRobotState} dumps full
  state to clipboard from the Telemetry GUI folder.
\end{itemize}

\subsection{Still to confirm during Phase 3 tuning
(resolved)}\label{still-to-confirm-during-phase-3-tuning-resolved}

\begin{itemize}
\tightlist
\item
  \st{Robot mass target} --- settled at \textbf{2.0 kg} with inertia
  \texttt{\{x:\ 0.014,\ y:\ 0.007,\ z:\ 0.009\}}.
\end{itemize}

\end{document}
